\section{Concluzii}
\label{sec:conclusion}

Lucrarea de față a analizat viabilitatea utilizării protocolului \textit{Push-Pull Dynamic Average Consensus} pentru monitorizarea sistemelor distribuite, ca alternativă la arhitecturile centralizate. Rezultatele obținute confirmă ipoteza că abordarea descentralizată oferă avantaje majore în termeni de scalabilitate și reziliență.

Prin comparația directă a topologiilor, am demonstrat că arhitectura \textbf{Mesh} oferă cel mai bun compromis pentru convergența rapidă ($\bigO{\log N}$), fiind superioară topologiilor Ring (prea lentă) și Star (vulnerabilă la congestie). De asemenea, testele de robustețe au evidențiat capacitatea algoritmului de a se adapta dinamic la schimbările de infrastructură. Deși pierderea bruscă a nodurilor cauzează o fluctuație a mediei estimate (\textit{Mass Loss}), sistemul nu intră în colaps, ci converge către o nouă stare stabilă.